% !TeX root = ../main.tex

\section{Выводы}

В ходе выполнения лабораторной работы было проведено комплексное исследование архитектуры и характеристик функциональных узлов квадратурного коррелятора аппаратуры потребителей СРНС ГЛОНАСС с использованием программно-аппаратной имитационной модели[cite: 1].

Практическое конфигурирование генератора дальномерного кода подтвердило полное соответствие алгоритма формирования псевдослучайной последовательности стандартной точности (СТ) требованиям интерфейсного контрольного документа ГЛОНАСС: использование 9-разрядного регистра сдвига с характеристическим полиномом $1 + x^5 + x^9$, единичной начальной инициализацией и съёмом сигнала с 7-го разряда обеспечивает генерацию периодического кода с периодом 511 чипов и длительностью $1\text{ мс}$[cite: 1, 2]. По спектральным характеристикам в полосе «Бесконечность» экспериментально подтверждены параметры сигнала: промежуточная частота составила $6{,}0\text{ МГц}$, а ширина полосы по первым спектральным нулям — $1{,}022\text{ МГц}$, что в точности соответствует удвоенной тактовой частоте дальномерного кода[cite: 1].

Исследование избирательного радиотракта показало, что сужение полосы пропускания фронтенда с $\infty$ до $6\text{ МГц}$ и $1\text{ МГц}$ ведёт к срезанию высших гармоник чипов ПСП и затягиванию их фронтов, что выражается в сглаживании острой вершины треугольной автокорреляционной функции и снижении крутизны дискриминатора кода[cite: 3, 5]. Одновременно наклон фазочастотной характеристики фильтра вносит аппаратурное групповое время запаздывания $\tau_{\text{гр}}$, составившее $0{,}573\text{ мкс}$ для полосы $6\text{ МГц}$ и $2{,}847\text{ мкс}$ для полосы $1\text{ МГц}$, что требует обязательной аппаратно-программной калибровки в приёмнике во избежание появления систематической ошибки псевдодальности[cite: 1].

Анализ шумовых и помеховых режимов продемонстрировал высокую эффективность когерентного сжатия спектра в корреляционном интеграторе[cite: 1]. Экспериментально оцененное отношение мощности сигнала к спектральной плотности шума $q_{c/n0} \approx 45{,}2\text{ дБГц}$ обеспечивает надежное обнаружение максимума АКФ даже в условиях, когда мгновенная мощность полезного сигнала на входе тракта многократно меньше мощности теплового шума[cite: 5].

Исследование однобитного квантования показало сохранение формы дискриминационной характеристики за счёт эффекта стохастического дизеринга при малых потерях отношения сигнал/шум (менее $2\text{ дБ}$), что делает применение компараторов технически и экономически оправданным при проектировании навигационных СБИС[cite: 3]. Моделирование гармонической помехи и пошаговая регистровая визуализация подтвердили высокую степень подавления сосредоточенных помех коррелятором и наглядно продемонстрировали механизмы разделения квадратур и формирования опережающей, точной и запаздывающей компонент для контуров слежения ФАП и ССЗ[cite: 1, 4, 5].